\section{Benchmark Construction: ToolPrivBench-SEA}
\label{sec:benchmark}

To evaluate cross-lingual over-privilege, we introduce \textsc{ToolPrivBench-SEA}, spanning three functional enterprise domains: database operations, filesystem access, and communication/broadcast services.

\subsection{Multilingual and Cross-Register Adaptation}
The benchmark features parallel task descriptions across three language tiers:
\begin{enumerate}
    \item \textbf{English (High-Resource Anchor):} Canonical formulations using direct syntactic imperatives.
    \item \textbf{Indonesian (Regional Lingua Franca):} Standard Indonesian phrasing reflective of national civic services.
    \item \textbf{Javanese (Low-Resource Regional Language):} Evaluated across both informal (\textit{Ngoko}) and polite/formal (\textit{Krama}) registers. The honorific nature of \textit{Krama} frequently uses passive nominalizations (e.g., \textit{``Nyuwun tulung dipuntingali...''}), providing an authentic test of model robustness against syntactic ambiguity.
\end{enumerate}

\subsection{Quality Assurance and Inter-Annotator Agreement}
All translations were reviewed by two independent bilingual annotators. We measured agreement using Cohen's Kappa:
\begin{equation}
    \kappa = \frac{P_o - P_e}{1 - P_e}
\end{equation}
Annotators achieved $\kappa = 0.88$, confirming strong inter-rater reliability across cultural and register adaptations.
